\chapter{Getting people back: routing and guidance}\label{ch:routing}

Once a casualty is located, two things can happen. If they can walk, Near guides them to the CCP by voice
and moves their pin as they walk. If they can't, Near plans the route for the medic team going out. Both
routes avoid the danger zones.

\section{The walking network}

\code{route.py} builds a graph from every walkable way in the cached OpenStreetMap data: residential streets,
footways, paths and tracks (\code{forest.geojson}), plus primary to tertiary roads (\code{lines.geojson}).
Motorways, trunk roads and tunnels are left out.
\begin{itemize}
  \item Points that share an OSM node have identical coordinates, so ways join there. Points closer than 4 m are
        also joined, which repairs small gaps in the drawing.
  \item Edges longer than 40 m are cut, so the danger check along them stays fine-grained.
  \item Result: 45,565 nodes and 75,768 edges, built in about 1 s. Start and end points snap to the largest
        connected component (44,781 nodes).
\end{itemize}

\section{The cost of a route}

Each edge $e$ of length $\ell_e$ costs
\begin{equation}
  c_e = \ell_e \left(1 + \sum_{z \in \text{zones}} w_{\text{type}(z)} \cdot \text{conf}_z \cdot
  \text{fresh}_z \cdot m_z(e)\right),
\end{equation}
where $m_z(e) \in [0,1]$ is the soft membership of the edge's midpoint in zone $z$ (1 inside, falling to 0 over
15 \% of the radius + 25 m), $\text{fresh}_z = 1 - \text{age}/\text{lifetime}$, and the type weights are
placeholders: mines 200, sniper 40, small arms 30, armour 25, explosions 15, drone 8. A weight of 30 means one
metre under small-arms fire is worth 30 extra metres of walking. Unplaced zones are ignored.

Dijkstra (scipy, sparse matrix) runs twice: once on this cost and once on plain length. The difference is what
the coordinator sees: ``+390 m detour, avoids 398 m inside the small-arms zone'' in the test scene. Each route
takes about 30 ms, and results are cached by start, end and risk map.

\nanofig{routing}{A walking route that bends around a reported small-arms zone, cut into spoken legs with landmark cues.}
{Top-down map illustration, 16:9, plain white background, flat vector, streets in light grey, buildings in pale blocks, a castle
keep icon. A start marker ``casualty'' at bottom-left and a ``CCP'' cross marker at top-right. A thin grey dashed straight-ish
route labelled ``shortest: 732 m'' passing through a translucent red circle labelled ``small-arms fire''. A thick cyan route
labelled ``safest: 1,122 m'' bending around the circle. Along the cyan route, numbered leg markers 1 to 5 at turns, with
small callouts: ``1: keep at your 7 o'clock, 30 m'', ``2: turn left, 40 m'', ``3: turn left, 100 m, keep on your right'',
``4: turn right'', ``5: CCP''. Exactly these labels. Colour palette (use only these): navy \#0D1110, cyan \#1AA7C4, amber \#E09A2B, red \#E5484D, green \#3FA877, grey \#8B9A94, plus white for text on dark shapes.}

\section{Spoken legs}

The route is simplified with Douglas--Peucker at 12 m, so wiggles smaller than that aren't announced as turns.
Nearly straight continuations (under 25$^\circ$) on the same street are merged, and so are legs under 15 m. Each leg
gets:
\begin{itemize}
  \item a turn word: straight on ($<$ 25$^\circ$), bear ($<$ 60$^\circ$), turn ($<$ 150$^\circ$), turn back;
  \item a length rounded for speech (10 m steps below 100 m, then 50 m);
  \item the street name if OSM has one;
  \item a landmark cue at its end: the nearest named landmark within 150 m and its side (left, right, ahead), not
        repeated from the previous leg;
  \item for soldiers who gave compass bearings, the leg's magnetic bearing, in degrees or in mils, matching the unit they used.
\end{itemize}

The first leg can't say ``turn left'', so it's oriented with what the speaker already gave us. That's their own
heading if it's known, otherwise a landmark they reported (``Walk with the Donjon at your 7 o'clock''), otherwise
the sun, otherwise a compass point. When a route changes mid-walk, the new first leg is given relative to the
current walking direction.

\section{Dead reckoning: the pin follows the walker}

The pin moves along the route by distance walked $s$, with an uncertainty radius $r$ that grows from the last
confirmation:
\begin{align}
  s &= s_0 + \begin{cases} 0.7\,\text{m} \times (\text{steps} - \text{steps}_0) & \text{phone step counter}\\
                            v \cdot \Delta t & \text{time only } (v = 1.1 \text{ m/s walking, } 0.6 \text{ assisted})\end{cases}\\
  r &= r_0 + k \cdot (s - s_0), \qquad k = 0.10 \text{ (steps)},\; 0.25 \text{ (time only)}.
\end{align}
\begin{itemize}
  \item With time only, the pin waits at the end of the current leg instead of guessing that the walker turned,
        and after 30 m of overshoot the phone says ``say next when you reach the turn''.
  \item ``Reached'' (button, or ``next'' said on the radio) snaps the pin to the junction and shrinks the radius to
        40 \% of what it was (at least 20 m).
  \item A pace count over the radio (``300 metres done'') sets the distance on the current leg directly.
  \item A new landmark report re-locates the walker around the dead-reckoning position (Chapter~\ref{ch:loc}), and the
        route restarts from there.
  \item ``I'm lost'' stops the guidance and asks for landmarks with clock directions.
\end{itemize}

\nanofig{deadreckoning}{The uncertainty of a walking casualty's position: it grows with distance walked and snaps back at each
confirmed turn.}
{Clean technical chart-style illustration, 16:9, plain white background. A horizontal axis ``distance walked (m)'' and a vertical
axis ``uncertainty radius (m)''. A cyan sawtooth line: rising slopes labelled ``+25\% of distance (time only)'' and a gentler
dashed amber line labelled ``+10\% (step counter)'', dropping sharply at three vertical markers labelled ``reached turn 1'',
``reached turn 2'', ``reached turn 3''. Above the chart, a small top-down strip of a walking route with a moving pin and a circle
that grows and shrinks. Minimal, exact labels only. Colour palette (use only these): navy \#0D1110, cyan \#1AA7C4, amber \#E09A2B, red \#E5484D, green \#3FA877, grey \#8B9A94, plus white for text on dark shapes.}

\section{Rerouting and warnings}

A background loop moves every guided walker every 2 s. When the set of danger zones changes, it checks the
rest of the walker's route.
\begin{itemize}
  \item If more than 20 m of the remaining route now lies inside a zone and a safer route exists, the walker gets
        the new route: ``New danger reported on your way. New route. Turn left, then about 80 metres.''
  \item If there's no safer route, the walker gets a spoken warning instead.
\end{itemize}
Danger zones are pinned where they were reported, so a walking reporter doesn't drag its own threats along.

\section{Medic route for casualties who can't walk}

For stretcher or trapped casualties located within 500 m, the map shows a dashed amber route from the CCP to the
casualty, with time out (1.3 m/s) and time back carrying a stretcher (0.6 m/s). Both speeds are placeholders.
The route uses the same risk-weighted cost, and is shown only once the position is no longer split between
several areas.

\section{Voice commands while walking}

``Next'' or ``reached'', ``repeat'', ``I'm lost'', ``arrived'', ``stop guidance'', in English and French
(\emph{suivant}, \emph{répète}, \emph{perdu}, \emph{arrivé}). A sentence longer than eight words is treated as a
report, not a command. The same commands come from the phone's buttons or, offline, inside a data packet.

\section{Why not reinforcement learning}

We were asked whether reinforcement learning (RL) could find better routes. For the route itself, no. Both the
network and the risk map are known when a route is needed, so the safest route is a shortest-path problem that
Dijkstra solves exactly in 30 ms. RL would learn an approximation of that answer, need a simulator and training,
and give a medic no reason they can check. Sequential decision-making under uncertainty does help in four places.
They're ranked below; the full write-up is in \code{reports/Route planning and reinforcement learning.md}.
\begin{enumerate}
  \item \textbf{The walker doesn't know exactly where he is.} Choose the instruction that works best on average over
        the position probabilities (a POMDP; the QMDP approximation uses a single cost-to-go map computed by one
        Dijkstra run from the CCP). About half a day of work, and our next step.
  \item \textbf{Threats found on the way, or moving.} This is the Canadian Traveller Problem. Our current policy,
        replanning on every new report, is the standard baseline; hedging policies do better.
  \item \textbf{Learning the risk weights} from after-action data with inverse RL. That needs data we don't have.
  \item \textbf{Dispatching several medic teams}, which is where military medical-evacuation research actually uses
        MDPs and RL (approximate dynamic programming, Q-learning).
\end{enumerate}

\begin{keyresult}[Tested (\code{tests/test\_guide.py}, 10 tests)]
Route avoids a small-arms zone placed on the direct path; simulated walker reaches the CCP with each turn confirmed and no
instruction spoken twice; clock-only mode waits at the turn; step counter moves the pin 120--160 m for 200 steps; a sniper
reported across the rest of the route triggers a reroute; voice commands; threats stay where reported; medic route for a
stretcher case; compass bearings in mils and pace count for a soldier without a phone.
\end{keyresult}
